\section{Keep caller data and private work separate}
The caller grants read access to two exact input slices and the four descriptor
words. It grants write access to the exact output slice. The backing allocations
are larger than many test matrices, so an access can remain mapped while violating
the logical boundary. A successful load does not prove that its address was allowed.

The checks vary each supplied base within its backing allocation. They inspect
full-width memory notices, including accesses that cross a selected boundary.
They also observe writes to input and descriptor storage. Writing a value and then
restoring it still violates read-only ownership, even if the final bytes look correct.

The tiled implementation reads partial output sums. Every byte it reads must first
have been initialized during the same call. Initial memory contents and a previous
call's product are not valid starting sums. A zero-reduction call is particularly
useful here: it must replace existing output contents with zeros while leaving both
inputs unread.

Use repeated calls in one execution session to test this lifetime. Compute a
nonempty product, then a zero-reduction product of another shape, then another
nonempty product. Reset all loop starts, shape values and partial sums at the
appropriate call boundary. A test harness that always creates a fresh memory image
can hide dependence on stale state.

\begin{MARSExercise}{Explain what a sentinel proves}
  Suppose the last output word still has its sentinel value after a call. Give two
  possible explanations, then describe the memory observation that distinguishes
  them. Explain why final sentinels alone cannot prove that an input was never written.
\end{MARSExercise}
